\section{Background}
\label{sec:background}

\subsection{Tokenization}

Tokenization is the process of converting raw text into a sequence of discrete
units (tokens) that a model can process. Two strategies are used in this project:

\begin{itemize}
    \item \textbf{Character-level tokenization.} Each character (letter, space,
          punctuation) is an independent token. The vocabulary is small
          ($\approx$60--80 symbols), but sequences are long and individual tokens
          carry little semantic meaning on their own.

    \item \textbf{Word-level tokenization.} Each whitespace-delimited word is a
          token. The vocabulary is larger but sequences are shorter. Rare or unseen
          words are mapped to a special \texttt{<UNK>} token.
\end{itemize}

\subsection{Markov Chain Language Models}

A Markov Chain model predicts the next token based on the previous $n$ tokens,
using frequencies observed in the training data. It stores a transition table
(a lookup of what came next after each n-gram) and samples from it during
generation. No neural network or gradient descent is involved — the model simply
counts and samples. It trains in milliseconds but cannot generalize beyond
patterns it has directly seen.

\subsection{Recurrent Neural Networks}

An RNN processes a sequence one token at a time, passing a hidden state forward
at each step. This hidden state acts as the model's memory of everything seen so
far. The network is trained by backpropagation through time (BPTT).

A key limitation is the \textbf{vanishing gradient problem}: as gradients flow
backwards through many timesteps, they shrink exponentially and become
uninformative. This prevents RNNs from learning dependencies between tokens that
are far apart in the sequence \cite{bengio1994}.

\subsection{Long Short-Term Memory}

LSTMs \cite{hochreiter1997} solve the vanishing gradient problem by introducing
a separate \textbf{cell state} and three \textbf{gates} that control what
information is kept, added, or discarded at each step:

\begin{itemize}
    \item \textbf{Forget gate} — decides what to erase from memory.
    \item \textbf{Input gate} — decides what new information to store.
    \item \textbf{Output gate} — decides what to output from the current memory.
\end{itemize}

Because the cell state is updated additively (not multiplicatively), gradients
can flow back over hundreds of timesteps without vanishing.

\subsection{Attention Mechanisms}

In a standard encoder-decoder model, the entire input sequence is compressed
into a single fixed-size vector before decoding begins. This creates a bottleneck
for longer inputs. Bahdanau attention \cite{bahdanau2015} addresses this by
allowing the decoder to look back at all encoder outputs at every step and decide
which parts of the input are most relevant at that moment. The result is a
weighted context vector that is fed into the decoder alongside the previous
predicted token.

\subsection{Transformer and Multi-Head Attention}

The Transformer \cite{vaswani2017} replaces recurrence entirely with
\textbf{self-attention}: every token in the sequence attends to every other token
simultaneously, in parallel. This removes the sequential bottleneck of RNNs and
allows the model to capture long-range relationships directly.

\textbf{Multi-head attention} runs several independent attention operations in
parallel, each learning to focus on different types of relationships (e.g.,
syntactic vs.\ semantic). Because there is no sequential processing, the model
cannot infer token order on its own — \textbf{sinusoidal positional encodings}
are added to the input embeddings to inject position information.
